\appendix
\huge{\textbf{Anexo}}
\section{Inductor}

\subsection{Características Adicionales}
Algunas características adicionales del inductor son:

\begin{itemize}
    \item Núcleo TDK RM10, material N30.
    \item 54 espiras de cable de $d = 0,3$ mm.
    \item $L_{gap} = 0,3$ mm
\end{itemize}

\subsection{Corriente máxima admisible}

La corriente máxima admisible del inductor está limitada por dos fenómenos: la saturación del núcleo y el límite térmico del hilo de cobre. La corriente máxima por saturación magnética resulta $I_{pico,sat} \approx 618$ mA.

Para un convertidor Boost, el campo magnético en el núcleo tiene una componente DC ($B_{DC}$) debida a la corriente media del inductor y una componente de ripple ($\Delta B$) debida a la conmutación. La condición límite es:

\begin{equation}
    B_{DC} + \frac{\Delta B}{2} \leq B_{max}
    \label{eq:blimite}
\end{equation}

donde $B_{max} = 130$ mT es el límite adoptado a $100°C$ para la ferrita N30.

El ripple generado por una cuadrada simétrica durante $t_{on} = D/f = 100\,\mu$s es:

\begin{equation}
    \Delta B = \frac{V_d \cdot t_{on}}{N \cdot A_{e,min}} = \frac{3,3 \cdot 100\times10^{-6}}{54 \cdot 90\times10^{-6}} \approx 67,9\text{ mT}
    \label{eq:deltaB}
\end{equation}

Despejando la componente DC máxima de (\ref{eq:blimite}):

\begin{equation}
    B_{DC,max} = B_{max} - \frac{\Delta B}{2} = 130 - 34 = 96\text{ mT}
    \label{eq:bdcmax}
\end{equation}

La corriente media máxima antes de saturar es:

\begin{equation}
    I_{medio,max} = \frac{B_{DC,max} \cdot N \cdot A_{e,min}}{L} = \frac{0,096 \cdot 54 \cdot 90\times10^{-6}}{1,02\times10^{-3}} \approx 457\text{ mA}
    \label{eq:imediomaxsat}
\end{equation}

El ripple de corriente en el inductor es:

\begin{equation}
    \Delta I_L = \frac{V_d \cdot D}{L \cdot f} = \frac{3,3 \cdot 0,5}{1,02\times10^{-3} \cdot 5000} \approx 323\text{ mA}
    \label{eq:deltaIL}
\end{equation}

Finalmente, la corriente pico máxima admisible por saturación es:

\begin{equation}
    I_{pico,sat} = I_{medio,max} + \frac{\Delta I_L}{2} = 457 + 161 \approx 618\text{ mA}
    \label{eq:ipicosat}
\end{equation}

\subsection{Límite térmico del hilo de cobre}

La sección transversal del hilo de $d = 0,3$ mm es:
\begin{equation}
    A_{Cu} = \pi \left(\frac{d}{2}\right)^2 = \pi \cdot (0,15)^2 \approx 0,071 \text{ mm}^2
    \label{eq:seccion}
\end{equation}

Adoptando $J_{max} = 3$ A/mm$^2$:
\begin{equation}
    I_{max,termico} = J_{max} \cdot A_{Cu} = 3 \cdot 0,071 \approx 213 \text{ mA}
    \label{eq:itermico}
\end{equation}

Dado que $I_{max,termico} < I_{pico,sat}$, el límite más restrictivo es el térmico, resultando en una corriente máxima admisible de $213$ mA. La fuente se limitó a $150$ mA como medida de seguridad.
